%% 10. CONCLUSION
%% ============================================================================

We have presented reasoning-core, a System 2 sidecar intended to address the
structural blind spot in LLM-based software development. By combining
Tree-sitter parsing for AST and call graph extraction with a Mamba SSM backbone
for semantic embedding, the system provides structural reasoning capabilities
that complement the linguistic strengths of LLMs.

The results below are historical controlled-study measurements from one
codebase, using an earlier risk-label schema. They are hypothesis-generating,
not evidence that the current default installation improves real coding
quality, reduces regressions, or lowers cost.

In that historical evaluation, 48 runs across 8 real-world engineering tasks
reported the following differences for the sidecar condition:

\begin{itemize}
    \item Reports +0.32 historical plan-quality and +0.20 implementation-quality differences (3-judge BARS scales)
    \item Reports 8.2 percent lower historical token usage on average (up to 29 percent on cache-heavy tasks)
    \item Observes 100 percent pass rates on locked and rotated tests for the sidecar arm (vs 92 percent/90 percent for vanilla)
    \item Reports higher historical diff discipline (+0.23)
    \item Reports higher historical repository-pattern fit (+0.43)
\end{itemize}

The system operates 100 percent locally with no external telemetry, addressing privacy concerns that have limited adoption of cloud-based code analysis tools. The sidecar added approximately 98 seconds per run in that study. Whether this
overhead is offset by fewer regression loops or lower cost in real usage
remains an open, outcome-linked question.

The evaluation methodology demonstrates the value of pre-registration,
kill-switches with documented amendment paths, and rigorous infrastructure.
The current next step is to validate those hypotheses against labeled existing
repository sessions.

The historical results motivate, but do not prove, the hypothesis that hybrid
System 1 + System 2 architectures can improve LLM-based software development.
LLMs excel at linguistic tasks, while specialized components may help with
structural reasoning; real-session outcome evidence is still required.

%% PLACEHOLDER: Add summary of contributions
%% PLACEHOLDER: Add impact statement
%% PLACEHOLDER: Add call to action for the research community
